\documentclass{article}
\usepackage{CJK}
\usepackage{amsmath,amssymb,graphicx,theorem,latexsym}

%%%%%%%%%% Start TeXmacs macros
\catcode`\<=\active \def<{
\fontencoding{T1}\selectfont\symbol{60}\fontencoding{\encodingdefault}}
{\theorembodyfont{\rmfamily}\newtheorem{definition}{Definition}}
\newtheorem{lemma}{Lemma}
\newcommand{\nonconverted}[1]{\mbox{}}
\newtheorem{theorem}{Theorem}
%%%%%%%%%% End TeXmacs macros

\begin{document}
\begin{CJK*}{UTF8}{gbsn}

\title{Introduction to Mathematical Logic --- Stynax}

\maketitle

{\tableofcontents}

\section{Argument论证}

\subsection{consist}

1.Declarative Sentence \ 陈述

2.Common Sence \ 常识

3.Inference Rules \ 推理

\

\section{Proposition}

\subsection{Proposition (命题)}

A proposition is a declarative sentence that can be judged as
{\underline{either true or false.}}

"This sentence is false."is not a proposition.

\subsection{Atomic Proposition (原子命题)}

A proposition that does not contain any smaller part that is still a
proposition.

\subsection{Compound Proposition (复合命题)}

A proposition that involves the {\underline{assembly}} of multiple
propositions is called a compound proposition.

\begin{tabular}{p{3.0cm}p{3.0cm}p{3.0cm}p{3.0cm}}
  Connectives & Read as & 名称 & Term\\
  \hline
  {\wedge} & and & 合取 & conjunction\\
  $\vee$ & or & 析取 & disjunction\\
  $\neg$ & not & 否定 & negation\\
  {\rightarrow} & if...then... & 蕴含 & implication\\
  {\leftrightarrow} & if and only if & 等价 & equivalence
\end{tabular}

\

\

\section{Syntax of propositional Logic
$\nonconverted{cal*L}^p$（命题逻辑的语法）}

Syntax--grammar

Semantics--language

\subsection{Alphabet of \  \nonconverted{cal*L} \textsuperscript{p}}

\begin{tabular}{|l|l|}
  \hline
  Atomic proposition & p,q,r,{...}, $p_1, q_1, r_2,${...}\\
  \hline
  Logical connectives & $\wedge, \vee, \neg, \rightarrow, \leftrightarrow$\\
  \hline
  Punctuation & ( , )\\
  \hline
\end{tabular}

\subsection{Expression of  \nonconverted{cal*L} \textsuperscript{p}}

Expression is the finite string of {\underline{symbols}}.

{\underline{The length of an expression is the number of occurrences of
symbols in it.}}

Two expressions u and v are {\underline{equal}} if they are of the same length
and have the same symbols in the same order.

\subsection{(Well-Formed) Formulas of  \nonconverted{cal*L}
\textsuperscript{p}}

\begin{definition}
  原子命题定义
  
  $\ensuremath{\operatorname{Atom}} (\nonconverted{cal*L}^p)$ is a set of
  expression of $\nonconverted{cal*L}^p
  \ensuremath{\operatorname{consisting}}\ensuremath{\operatorname{of}}$ an
  atom proposition symbol only.
\end{definition}

\begin{definition}
  合式公式的定义Well Formed Formula
  
  $\ensuremath{\operatorname{Form}} (\nonconverted{cal*L}^p)$ is the smallest
  set of expression that satisfies (i)(ii)(iii)
  
  (i) $\ensuremath{\operatorname{Atom}} (\nonconverted{cal*L}^p) \subseteq
  \ensuremath{\operatorname{Form}} (\nonconverted{cal*L}^p)$
  
  (ii)If $\alpha \in \ensuremath{\operatorname{Form}}
  (\nonconverted{cal*L}^p), \ensuremath{\operatorname{then}} (\neg \alpha) \in
  \ensuremath{\operatorname{Form}} (\nonconverted{cal*L}^p) .$
  
  (iii)If $\alpha, \beta \in \ensuremath{\operatorname{Form}}
  (\nonconverted{cal*L}^p), \ensuremath{\operatorname{then}} (a \wedge \beta),
  (\alpha \vee \beta), (\alpha \rightarrow \beta), (\alpha \leftrightarrow
  \beta) \in \ensuremath{\operatorname{Form}} (\nonconverted{cal*L}^p)$
  
  注意，这种定义的写法存在一定问题：它只说明了存在条件，并没有说明条件唯一。
  
  事实上，合式公式的定义基于递归结构，只能是由(i)(ii)(iii)生成的有限表达式。
\end{definition}

\subsection{Common Properties of Well-Formed Formula}

\begin{lemma}
  Well-formed formula of $\nonconverted{cal*L}^p$ are non-empty
  expression.（合式公式非空定理）
\end{lemma}

\begin{lemma}
  Every well-formed formula of $\nonconverted{cal*L}^p$ has equal number of
  openning and closing bracket.
  
  （合式公式括号匹配定理）
\end{lemma}

\noindent\textbf{Proof\ }We will prove "$\ensuremath{\operatorname{Form}}
(\nonconverted{cal*L}^p)$ of all length has equal number of openning and
closing bracket" by induction. Let n=the length of the formula.
\begin{itemize}
  \renewcommand{\labelitemi}{$\bullet$}\renewcommand{\labelitemii}{$\bullet$}\renewcommand{\labelitemiii}{$\bullet$}\renewcommand{\labelitemiv}{$\bullet$}\item
  Base Case :n=1:
  
  For p is the WFF of length 1.
  
  p $\in$ Atom($\nonconverted{cal*L}^p$), is the simpliest WFF.
  
  $\#_{(} (p) =\#_{)} (p) = 0 $
  
  \item Induction Hypothesis :n<k:
  
  Suppose all WFF of length less than k has equal number of openning and
  closing bracket.
  
  \item Induction Step: Now we consider n=k.
  
  Let p be a WFF of length k.
  
  Case 1:
  
  p=$(\neg q)$, then $\#_{(} (p) =\#_{(} (q) +
  1\ensuremath{\operatorname{and}} \#_{)} (p) =\#_{)} (q) + 1$
  
  Because q is a WFF of length k-3, less than k,
  
  then according to H.I. $\#_{(} (q) =\#_{)} (q)$
  
  So $\#_{(} (p) =\#_{)} (p)$
  
  Case 2:
  
  p=($q_1 \ensuremath{\operatorname{op}}q_2$) in which op$\in \{ \wedge, \vee,
  \rightarrow, \leftrightarrow \}$
  
  then $\#_{(} (p) =\#_{(} (q_1) +\#_{(} (q_2) + 1$ and $\#_{)} (q) =\#_{)}
  (q_1) +\#_{)} (q_2) + 1$
  
  Because \#$(q_1) \leqslant k - 4 < k,
  \ensuremath{\operatorname{so}}\ensuremath{\operatorname{does}}\# (q_2)$
  
  then according to H.I. \ $\#_{(} (q_1) =\#_{)} (q_1)
  \ensuremath{\operatorname{and}}\#_{(} (q_2) =\#_{)} (q_2)$
  
  So $\#_{(} (p) =\#_{)} (p)$
  
  \item Summary:
  
  WFF of all length has equal number of openning and closing bracket.
  
  \ 
\end{itemize}
\hspace*{\fill}$\Box$\medskip

\begin{otherlanguage}{chinese}
  \begin{lemma}
    (真前后缀括号不完备定理)
    \begin{itemize}
      \renewcommand{\labelitemi}{$\bullet$}\renewcommand{\labelitemii}{$\bullet$}\renewcommand{\labelitemiii}{$\bullet$}\renewcommand{\labelitemiv}{$\bullet$}\item
      Every proper prefix（真前缀） of a well formed formula in
      $\nonconverted{cal*L}^p$ has more openning bracket than closing bracket.
      
      \item Every proper suffix（真后缀） of a well-formed formula in
      $\nonconverted{cal*L}^p$ has more closing bracket than openning bracket.
      
      \item Hence, proper prefix and proper suffix are not wff in
      $\nonconverted{cal*L}^p$.
    \end{itemize}
  \end{lemma}
\end{otherlanguage}

\noindent\textbf{Proof\ }Let's prove the theorem about prefix by structural
induction, that about suffix has a similar proof:

Let $P (\varphi)
\ensuremath{\operatorname{be}}\ensuremath{\operatorname{the}}\ensuremath{\operatorname{property}}\ensuremath{\operatorname{that}}$
every proper prefix of the well-formed formula {\varphi} has more opening
brackets than closing brackets.
\begin{itemize}
  \renewcommand{\labelitemi}{$\bullet$}\renewcommand{\labelitemii}{$\bullet$}\renewcommand{\labelitemiii}{$\bullet$}\renewcommand{\labelitemiv}{$\bullet$}\item
  Base Case : If $\alpha$ is an atom, it has no proper prefixes and $P
  (\alpha)$ is true.
  
  \item Inductive Hypothesis : suppose $\varphi_1, \varphi_2$ be WFF
  satisfying $P (\varphi_1), P (\varphi_2)$.
  
  \item Inductive Step :
  
  \ \ \ \ \ Case 1:
  
  $p = (\neg \varphi),$
  
  let $\varphi'$ be any proper prefix of $\varphi$, so does $p$,then the
  proper prefix of $p$ can be:
  
  1.( ,$P (() $is true.
  
  2.($\neg$ , $P ((\neg ) $is true.
  
  3.($\neg \varphi'$ , $\#_{(} (( \neg \varphi') =\#_{(} (\varphi') + 1,
  \#_{)} (( \neg \varphi') =\#_{)} (\varphi')$
  
  Because $P (\varphi)$ is true. $\#_{(} (\varphi') >\#_{)} (\varphi')$
  
  so $\#_{(} (( \neg \varphi') >\#_{)} (( \neg \varphi') + 1 >\#_{)} (( \neg
  \varphi')$
  
  4.$(\neg \varphi , \#_{(} (\varphi) =\#_{)} (\varphi)
  \ensuremath{\operatorname{by}}\ensuremath{\operatorname{lamma}}2$, so it is
  obvious that $P ((\neg \varphi )
  \ensuremath{\operatorname{is}}\ensuremath{\operatorname{true}}.$
  
  \ \ \ \ \ Case2 :
  
  $p = (\varphi_1 \ensuremath{\operatorname{op}} \varphi_2),
  \ensuremath{\operatorname{in}}\ensuremath{\operatorname{which}}\ensuremath{\operatorname{op}}
  \in \{ \wedge, \vee, \rightarrow, \leftrightarrow \}$,
  
  proper prefix of $p$ can be:
  
  $\left(  \quad, \quad  \right( \varphi_1' \hspace{1em}, \quad \left(
  \varphi_1 \quad, \quad \right( \varphi_1 \ensuremath{\operatorname{op}}
  \quad, \quad \left( \varphi_1 \ensuremath{\operatorname{op}} \varphi_2'
  \quad \right. \ensuremath{\operatorname{and}} \quad (\varphi_1
  \ensuremath{\operatorname{op}} \varphi_2 $
  
  We can easily prove that all of them satisfy $P (p')
  \ensuremath{\operatorname{is}}\ensuremath{\operatorname{true}}.$
  
  \item Summary
  
  Every proper prefix of the well-formed formula p satisfies $P (p')$ is true.
\end{itemize}
\hspace*{\fill}$\Box$\medskip

\subsection{Parse Tree of $\nonconverted{cal*L}^p$}

\begin{definition}
  解析树定义
  
  A parse tree of a formula in $\nonconverted{cal*L}^p
  \ensuremath{\operatorname{is}}a$tree that:
  \begin{itemize}
    \renewcommand{\labelitemi}{$\bullet$}\renewcommand{\labelitemii}{$\bullet$}\renewcommand{\labelitemiii}{$\bullet$}\renewcommand{\labelitemiv}{$\bullet$}\item
    The root is a formula.
    
    \item Leaves are atoms.
    
    \item Each internal node is formed by applying some formation rule on its
    children.
  \end{itemize}
\end{definition}

\begin{lemma}
  解析树和合式公式的等价性
  
  An expression of $\nonconverted{cal*L}^p \ensuremath{\operatorname{is}}$a
  well-formed formula if and only if there is a parse tree of it.
\end{lemma}

$((p \vee q) \rightarrow ((\neg p) \leftrightarrow (q \wedge r)))$

\begin{align}
  \includegraphics[width=4.159238488783943cm,height=2.996130132493769cm]{PL1：Syntax-1.pdf}
  &  \nonumber\\
  \ensuremath{\operatorname{Simplified}}\ensuremath{\operatorname{Parse}}\ensuremath{\operatorname{Tree}}:
  &  \nonumber\\
  \hspace{-0.016791289518562246cm}\includegraphics[width=1.8328987931260659cm,height=2.8848386462022826cm]{PL1：Syntax-2.pdf}\hspace{-0.0041978223796406144cm}
  &  \nonumber
\end{align}

\begin{definition}
  子公式定义
  
  A formula G is a subformula of formula F if G occurs within F. G is a proper
  subformula of F if G $\neq$ F.
  
  The nodes of the (unsimplified) parse tree of F form the set of subformulas
  of F.
\end{definition}

\subsection{The Unique Readability Theorem}

\begin{theorem}
  唯一可读性定理
  
  There is a unique way to construct every well-formed formula.
  
  这条定理指出了，任意一个合式公式，只能唯一地分解成一个个最小的命题符号，以及唯一确定它们构建该合式公式的方式。
\end{theorem}

这是由expression的等价性决定的，``Two expressions u and v are
{\underline{equal}} if they are of the same length and have the same symbols
in the same order.''

\noindent\textbf{Proof\ }\

\begin{align}
  \varphi = \left\{\begin{array}{l}
    (\neg \alpha)  \left\{\begin{array}{l}
      \varphi = (\neg \alpha') \quad \Rightarrow \alpha = \alpha'\\
      \varphi = (\alpha' \ast \beta') \Rightarrow \alpha'
      \ensuremath{\operatorname{starts}} \neg \Rightarrow
      \ensuremath{\operatorname{contradictory}}
    \end{array}\right.\\
    \\
    (\alpha \ast \beta) \left\{\begin{array}{l}
      \varphi = (\neg \alpha') \Rightarrow \alpha
      \ensuremath{\operatorname{starts}} \neg \Rightarrow
      \ensuremath{\operatorname{contradictory}}\\
      \\
      \\
      \varphi = (\alpha' \ast \beta') \left\{\begin{array}{l}
        a, b = a', b'\\
        \alpha
        \ensuremath{\operatorname{is}}a\ensuremath{\operatorname{profer}}\ensuremath{\operatorname{prefix}}\ensuremath{\operatorname{of}}
        \alpha' \ensuremath{\operatorname{or}}\ensuremath{\operatorname{vise}}
      \end{array}\right.
    \end{array}\right.
  \end{array}\right. &  \nonumber
\end{align}

\ \hspace*{\fill}$\Box$\medskip

\

现观察异或公式：

\begin{align}
  (p \vee q) \wedge (\neg (p \wedge q)) = (p \wedge (\neg q)) \vee (q \wedge
  (\neg p)) & =\ensuremath{\operatorname{XOR}} (p, q) \nonumber
\end{align}

这两个公式逻辑等价，但他们是不同的合式公式，不违背唯一可读性定理。

\

\section{Conventions}

For readability, we use the following convention, which states that we may
drop the outermost parentheses.

\subsection{Parentheses}

Parentheses are used to resolve ambiguity.

If F or (F) is a formula, then we view F and (F) as the same formula

\subsection{Precedence and Associativity}

If no parentheses are present, we could use precedence and associativity to
disambiguate formulas.
\begin{itemize}
  \renewcommand{\labelitemi}{$\bullet$}\renewcommand{\labelitemii}{$\bullet$}\renewcommand{\labelitemiii}{$\bullet$}\renewcommand{\labelitemiv}{$\bullet$}\item
  Priority:$() \hspace{1em} > \quad \neg \quad > \quad \wedge \quad > \quad
  \hspace{1em} \vee \quad > \quad \rightarrow \quad > \quad \leftrightarrow$
  
  \item Associtivity:
  
  $\wedge, \vee \ensuremath{\operatorname{and}} \leftrightarrow
  \ensuremath{\operatorname{are}}\ensuremath{\operatorname{associative}}
  (阅读时默认按从左到右的自然顺序，若顺序相同可省略括号)$
  
  {\rightarrow} is right associative($默认按从右到左的顺序$)
\end{itemize}
\begin{align}
  \neg p \rightarrow p \wedge \neg q \vee r \leftrightarrow q &  \nonumber\\
  (\neg p) \rightarrow p \wedge (\neg q) \vee r \leftrightarrow q & 
  \nonumber\\
  (\neg p) \rightarrow (p \wedge (\neg q) \vee r) \leftrightarrow q & 
  \nonumber\\
  (((\neg p) \rightarrow (p \wedge (\neg q) \vee r)) \leftrightarrow q) & 
  \nonumber
\end{align}

\

\

\

\

\end{CJK*}
\end{document}
